
\section{概形可视化：一般点}

\epigraph
  {从前，人们常郑重其事地警告说：图画不严格；这一吓唬从未被人叫破，
   倒把它的受害者们永远地吓住，使其一味只求稳当。诚然，有些图画并不严格，
   但我得说，大多数图画是严格的（而我自己但凡可行，总要使用它们）。}
  {——J.~E.~\trLittlewood，[Lit, p.~54]}
\epigraph
  {我们都知道，艺术不是真理。艺术是一种谎言，它使我们领悟真理——至少是
   那被赋予我们去理解的真理。艺术家必须懂得以何种方式使人相信他的谎言的
   真实性。}
  {——P.~\trPicasso，[Pi, p.~315]}

多年来，你一直能够在平面中想象 $x^2+y^2=1$，而如今你对如何想象
$\Spec \mathbb{Z}$ 也已有了一些想法。既然我们宣称要（借助 $\Spec$
\trFunctor）把\trRing{}当作几何对象来理解，那么我们就应当希望对这些对象
发展出几何上的洞见。要培养关于\trScheme{}的几何直观，在头脑中备有图画是
有益的——把你对已然熟悉的那些\trGeometricSpace{}所怀的直观延伸过来便是。
随着行文推进，我们会凭经验逐步形成「\trScheme{}应当长成什么模样」的
一些想法。本节便总结我们迄今之所获。

有些数学家宁愿全然代数式地思考，从不以图画的方式想问题。另一些人则会因
这乃是艺术而非科学而心生不安。最后，这类「挥手比划」（hand-waving）式的
议论\footnote{译注：hand-waving，数学行话，指以挥挥手带过细节、不作严格
论证的讲法。}注定绝不会用于理论的严格展开。鉴于这些理由，你或许想跳过
这几节。然而，头脑中若有正确的图画，会大大有助于理解哪些事实应当为真、
又当如何证明它们。\trFitzGerald{}关于拓展心智之重要的那番勉励，是给数学家的
一条要紧忠告。

\epigraph
  {检验一流智力的标准，是看你能否在头脑中同时容纳两种相反的观念，
   而仍不丧失行事的能力。}
  {——F.~Scott~\trFitzGerald，《崩溃》[Fi, p.~41]}

我们的出发点是「仿射复\trVariety」这一例子（即在 $\mathbb{C}$ 上由涉及
有限多个变量的方程切出来的东西），以及更一般地，任意\trAlgClosedField{}
上的类似例子。我们从直观的观念入手（「传统的点」按你所期待的方式行事），
然后再加入两个既新颖又令人不安的特征：\trGenericPoint{}与非\trReduced{}行为。
此后你便可把这个观念拓展到那些貌似不同的空间，例如 $\Spec \mathbb{Z}$。

\trHilbert{}弱\trNullstellensatz（\cref{thm:hilbert-weak-nullstellensatz}）表明，
「传统的点」确实作为\trScheme{}的点而存在，且这一点对任何\trAlgClosedField{}皆然。
若域不是代数闭的，传统的点就会经\trGalois{}\trConjugate{}被粘成一团一团，
正如 \S3.2 中的例~5（实\trAffineLine）与例~6（$\mathbb{F}_p$ 上的%
\trAffineLine）那样。这是\trHilbert{}\trNullstellensatz（\cref{thm:hilbert-nullstellensatz}）
的一个几何解读。

但图中还有一些额外的点有待添入。你应当记住：这些点「对应」于「\trIrreducible」
「闭」的（代数）子集。作为动机，考虑复仿射平面（例~7）的情形：对每个%
\trIrreducible{}多项式，我们都曾有一个这样的点；此外另有一个点对应于整个平面。
当我们（在下一节中）定义\trZariskiTopology{}时，会把「闭」说得精确。
对于「\trIrreducible」应当是什么意思，你或许已有想法；我们将在
\S3.6 的开头把它说精确。所谓「对应」，意思是：每个闭的\trIrreducible{}子集都有
一个相应的点坐在它上面，称为它的\trGenericPoint（定义见 \S3.6）。它是一个新点，与该子集中所有其余的点都不相同。（对
$\Spec A$ 而言，这一对应描述于练习~3.7.F；对一般的\trScheme{}，则描述于
练习~5.1.B。）我们并不知道该把一般点画在何处才算精确，故不妨把它随意搁
在任何地方；但你应当把它想成「几乎处处」都在——特别地，它靠近该子集中
每一个其他的点。

在 \S3.2.8 中，我们曾看到 $\Spec A/I$ 的点应当如何解读为 $\Spec A$ 的
一个子集。所以举例来说，当你见到 $\Spec \mathbb{C}[x,y]/(x+y)$ 时，你不应
只把它想象成一条直线，而应想象成 $xy$ 平面中的一条直线：\trAlgebra{}
$\mathbb{C}[x,y]$ 的生成元 $x$ 与 $y$ 的选取，本身便蕴含了一个到仿射空间
中的\trInclusionMap{}。

在 \S3.2.9 中，我们曾看到 $\Spec S^{-1}A$ 的点应当如何解读为 $\Spec A$ 的
子集。两种最重要的情形都已在彼处讨论过：$\Spec A_f$ 的点对应于 $\Spec A$
中使 $f$ 不消失的点；我们稍后（\S3.5）将把它解读为一个\trDistinguishedOpen。

若 $\mathfrak{p}$ 是一个\trPrimeIdeal{}，则应把 $\Spec A_{\mathfrak{p}}$
看作「空间 $\Spec A$ 在 $\mathfrak{p}$ 所对应子集近旁的一片碎片」。此事
最简单的非平凡情形是 $A=k[x]$ 而 $\mathfrak{p}=(x)\subset A$（见
\cref{ex:small-schemes-dual-numbers}；我们将在练习~3.4.K 中再次讨论它）。

\epigraph
  {「他们中间若有谁能解释它，」爱丽丝说（她在方才这几分钟里已长得
   如此高大，打断他半点也不再害怕），「我就给他六便士。我才不信这里面有
   一丁点儿意思。」

   「这里面若是没有意思，」国王说，「那倒省去了天大的麻烦，你知道的，
   咱们也就不必费心去找了。」}
  {——刘易斯·\trCarroll，[Carr, Ch.~XII]}
